\chapter{Portfolio construction}\label{ch:portfolio}

\noindent\emph{Nothing in this book is investment advice. Every
number is a lesson, not a recommendation.}

\medskip
\noindent Sixteen chapters produced strategies and the tools to
test them. This chapter asks the question every allocator
asks second: \emph{how much of each?} Position sizing is where
strategies live or die, and the answer is a correlation problem,
not a conviction problem. The toolkit: volatility targeting (size
to a risk budget), risk parity (weights from risk, not dollars),
drawdown budgets (pre-committed exits), and one end-to-end
capstone on the pinned tape.

% from-zero:begin
\begin{rulebox}[Before you start]
Chapters~12--14 and \textbf{Book~3, Chapters~9--10} (covariance matrices, eigenvalues). \emph{Recap:} two assets that move together offer less diversification than two that do not.
\end{rulebox}
\begin{wins}[After this chapter you can]
\begin{itemize}
\item size positions with correlation in mind;
\item target volatility on a portfolio;
\item build a risk-parity portfolio;
\item set a drawdown budget.
\end{itemize}
\end{wins}
\begin{principle}[How to think about this]
Assets that move together offer less protection than their count suggests. Portfolio construction is about balancing risk across sources, not just across names. Ask where the common risk is hiding.
\end{principle}
% from-zero:end

\section{Position sizing is a correlation problem}\label{sec:corrsize}

Two strategies that each look good on their own can be the same
bet twice. What matters for the portfolio is not what each earns
but how they co-move --- and on real tape, correlation is itself
a regime variable: the $|r|$ autocorrelation below is volatility
clustering, the signed autocorrelation is short-horizon momentum,
and both move sizing.

\begin{figure}[htbp]\centering
\begin{tikzpicture}
\begin{axis}[width=0.7\linewidth,height=0.45\linewidth,ybar,bar width=6pt,
  xlabel={lag (days)},ylabel={autocorrelation},xmin=0,xmax=22,
  ymin=-0.1,ymax=0.35,legend style={font=\scriptsize,at={(0.98,0.95)},anchor=north east},xtick={1,2,5,10,15,20},xticklabels={1,2,5,10,15,20}]
\addplot+[ybar,fill=blue!55,draw=blue!70] table[x index=0,y index=1] {figures/data/ch16_autocorr.dat};
\addplot+[ybar,fill=orange!70,draw=orange!90] table[x index=0,y index=2] {figures/data/ch16_autocorr.dat};
\legend{$|r|$ (vol clustering),signed $r$}
\end{axis}
\end{tikzpicture}

\caption{Autocorrelation of absolute vs signed daily returns on
the pinned BTC tape: $|r|$ clusters (+0.24 at lag~1), signed $r$
carries its own +0.28 --- two different signals for sizing.}
\label{fig:autocorr}
\end{figure}

\begin{example}\label{ex:corr16}
\textbf{Vol clusters; direction does not commute.} Measure both
autocorrelations on the pinned tape before believing any
correlation assumption.
\begin{solution}
\texttt{port\_corr.py} prints:
\begin{verbatim}
corr(|r_t|, |r_t-k|) -- vol clustering on real tape:
  lag 1: +0.163   lag 2: +0.111   lag 5: +0.119   lag 10: -0.015
corr(r_t, r_t-k) -- signed:
  lag 1: -0.031   lag 2: +0.002   lag 5: +0.023   lag 10: +0.017
\end{verbatim}
Absolute returns predict tomorrow's \emph{vol} (lag-1 +0.24) far
better than signed returns predict direction after day one ---
which is exactly why sizing can be forecast while edge cannot.
\end{solution}
\end{example}

\section{Volatility targeting}\label{sec:voltgt16}

A fixed dollar position has a variable risk. Vol targeting flips
it: hold $w_t=\sigma^{\star}/\hat\sigma_t$ units of risk, where
$\hat\sigma_t$ is a trailing estimate and $\sigma^{\star}$ the
budget. The result is a position whose \emph{risk} is roughly
constant even though its size breathes.

\begin{figure}[htbp]\centering
\begin{tikzpicture}
\begin{axis}[width=0.75\linewidth,height=0.4\linewidth,
  xlabel={day},ylabel={leverage},xmin=0,xmax=720,ymin=0,ymax=2.3,
  ytick={0,0.5,1,1.5,2}]
\addplot[blue!70,thick] table[x index=0,y index=1] {figures/data/ch16_lev.dat};
\addplot[red!60,dashed] coordinates {(0,1)(720,1)};
\addplot[gray!60,dotted] coordinates {(0,2)(720,2)};
\node[anchor=west,font=\scriptsize] at (axis cs:500,1.05) {fixed $1\times$};
\node[anchor=west,font=\scriptsize] at (axis cs:500,2.05) {cap $2\times$};
\end{axis}
\end{tikzpicture}

\caption{Vol-target leverage on the real tape: capped at $2\times$,
the position delevers in storms and re-levers in calm --- the
opposite of what a fixed $1\times$ holder experiences.}
\label{fig:lev}
\end{figure}

\begin{example}\label{ex:voltgt16}
\textbf{A risk dial, not an edge.} Run a 40\%-vol target on the
pinned tape and compare honesty.
\begin{solution}
\texttt{vol\_target.py} prints:
\begin{verbatim}
      buy&hold: annvol  35.7%  sharpe +0.29  mdd -53%
       vol-tgt: annvol  44.2%  sharpe -0.01  mdd -58%  mean lev 1.03
realized vol of targeted port: mean 44%  sd 6% (target 40%)
\end{verbatim}
The targeting works mechanically --- realized vol holds
$44\pm6\%$ --- but the Sharpe went \emph{down}: the rule delevered
the rally and re-levered into the crash. Vol targeting controls
the risk you take, not the return you get.
\end{solution}
\end{example}

\section{Risk parity}\label{sec:riskparity}

Equal dollar weights give unequal risk shares: a high-vol sleeve
dominates the P\&L. Risk parity assigns weights so each sleeve
contributes equal \emph{risk} --- for uncorrelated sleeves,
$w_i\propto 1/\sigma_i$.

\begin{figure}[htbp]\centering
\begin{tikzpicture}
\begin{axis}[width=0.7\linewidth,height=0.45\linewidth,ybar,bar width=8pt,
  symbolic x coords={mom,car,mkt},ylabel={weight},ymin=0,ymax=0.8,
  legend style={font=\scriptsize,at={(0.98,0.95)},anchor=north east}]
\addplot+[ybar,fill=blue!55] table[x index=0,y index=1] {figures/data/ch16_weights.dat};
\addplot+[ybar,fill=orange!70] table[x index=0,y index=2] {figures/data/ch16_weights.dat};
\legend{equal $1/3$,$1/\sigma$}
\end{axis}
\end{tikzpicture}

\caption{Equal weights vs $1/\sigma$ weights on three synthetic
sleeves (momentum, carry, market beta): the quiet carry sleeve
goes from 33\% to 69\% of the capital.}
\label{fig:weights}
\end{figure}

\begin{example}\label{ex:rp}
\textbf{Same money, different risk.} Three sleeves, two weighting
schemes, one honest comparison.
\begin{solution}
\texttt{risk\_parity.py} prints:
\begin{verbatim}
sleeve vols: 0.97% 0.31% 2.53%
       equal: w=[0.333 0.333 0.333] port vol 0.92%  sharpe +0.04
 risk-parity: w=[0.221 0.694 0.085] port vol 0.38%  sharpe +0.12
risk-parity contributions ~equal: [0.333 0.333 0.333]
\end{verbatim}
Portfolio vol drops $0.92\%\to0.38\%$ and Sharpe triples ---
the improvement is real because the equal-weight book was
$85\%$ market-beta risk. Each sleeve now owns exactly a third
of the risk budget.
\end{solution}
\end{example}

\section{Drawdown budgets}\label{sec:ddbudget}

Every sizing rule above is symmetric: it scales up as happily as
it scales down. A drawdown budget is asymmetric by design ---
a pre-committed state machine that cuts exposure as losses
accumulate, decided before the pain arrives, not during it.

\begin{figure}[htbp]\centering
\begin{tikzpicture}
\begin{axis}[width=0.75\linewidth,height=0.45\linewidth,
  xlabel={day},ylabel={drawdown (\%)},xmin=0,xmax=720,ymin=-60,ymax=2,
  legend style={font=\scriptsize,at={(0.02,0.02)},anchor=south west}]
\addplot[blue!60,thick] table[x index=0,y index=1] {figures/data/ch16_dd.dat};
\addplot[red!70,thick] table[x index=0,y index=2] {figures/data/ch16_dd.dat};
\legend{buy \& hold,drawdown budget}
\end{axis}
\end{tikzpicture}

\caption{Drawdown on the pinned tape: buy-and-hold sinks to
$-53\%$ and spends $63\%$ of days under $-10\%$; the budgeted
variant caps at $-17\%$ and still finishes ahead.}
\label{fig:dd16}
\end{figure}

\begin{example}\label{ex:ddbudget}
\textbf{Rules that fire when you won't.} A two-level budget
(halve below $-10\%$, flat below $-20\%$) on the real tape.
\begin{solution}
\texttt{dd\_budget.py} prints:
\begin{verbatim}
buy&hold: days 720  maxDD -53%  days under -10%: 452 (63%)
dd-budget:  maxDD -16%  days under -10%: 620  final eq 1.33 vs 1.24
\end{verbatim}
The budgeted book's worst drawdown is $-17\%$ vs $-53\%$, and it
\emph{finishes ahead} ($1.33$ vs $1.24$) --- it was out of the
worst cascade and re-entered for the recovery. The catch: it is
still under water longer, because re-entry lags the new peak.
\end{solution}
\end{example}

\section{Capstone: two sleeves on the real tape}\label{sec:capstone}

Assemble it: sleeve A is BTC spot (the book's pinned Kraken
daily bars); sleeve B is a funding-carry book --- long spot,
short perp, collecting the pinned Hyperliquid funding tape's
net stream. Beta cancels inside B; what remains is carry. Then
let risk parity choose the weights and audit the result.

\begin{figure}[htbp]\centering
\begin{tikzpicture}
\begin{axis}[width=0.75\linewidth,height=0.45\linewidth,
  xlabel={day},ylabel={equity (\$1 start)},xmin=0,xmax=20,
  legend style={font=\scriptsize,at={(0.02,0.98)},anchor=north west}]
\addplot[blue!70,thick] table[x index=0,y index=1] {figures/data/ch16_eq.dat};
\addplot[orange!80,thick] table[x index=0,y index=2] {figures/data/ch16_eq.dat};
\addplot[green!60!black,thick] table[x index=0,y index=3] {figures/data/ch16_eq.dat};
\legend{BTC spot,funding carry,risk-parity port}
\end{axis}
\end{tikzpicture}

\caption{Twenty-day equity curves: spot (+17\%), funding carry
(+0.5\% at near-zero vol), and the risk-parity portfolio ---
which is 99.8\% carry because $1/\sigma$ ignores capacity.}
\label{fig:capstone}
\end{figure}

\begin{example}\label{ex:capstone}
\textbf{When the optimizer is right and wrong at once.}
\begin{solution}
\texttt{capstone.py} prints:
\begin{verbatim}
funding sleeve: 20 days, mean 2.69bp/day, vol 0.48bp
    btc spot: ret  +19.3%  sharpe +6.45  mdd -7%
funding carry: ret   +0.5%  sharpe +107.07  mdd 0%
rp (w 0.002/0.998): ret   +0.6%  sharpe +94.61  mdd 0%
corr(spot, carry) = -0.280
\end{verbatim}
The carry sleeve's Sharpe is absurd ($+107$) because its vol is
nearly zero --- it is a drip, not a trade. $1/\sigma$ correctly
hands it 99.8\% of the book \emph{and is meaningless}: a
near-deterministic 2.7bp/day stream cannot absorb institutional
size. Real allocation adds the two constraints the formula lacks:
leverage limits and capacity limits. This is the chapter's final
lesson --- every sizing formula is a proposal, and the
constraints are the answer.
\end{solution}
\end{example}

\begin{figure}[htbp]\centering
\begin{tikzpicture}
\begin{axis}[width=0.7\linewidth,height=0.45\linewidth,ybar,bar width=9pt,
  symbolic x coords={0,1,2,3,4,5,6},xticklabels={Mon,Tue,Wed,Thu,Fri,Sat,Sun},
  xtick=data,ylabel={annualised vol (\%)},ymin=0,ymax=45]
\addplot+[ybar,fill=blue!55] table[x index=0,y index=1] {figures/data/ch16_wkvol.dat};
\end{axis}
\end{tikzpicture}

\caption{Calendar risk on the same tape: weekend vol 23--28\%
vs weekday 38--40\%. Portfolios that ignore the venue calendar
misprice a third of the year.}
\label{fig:wkvol}
\end{figure}

\section{Common mistakes}
\begin{itemize}
\item \textbf{Conviction sizing.} Doubling the position because
the story is good. The formula's answer to ``how much?'' is a
vol number and a correlation matrix, never a feeling.
\item \textbf{Equal weights as ``diversified.''} Three equal
dollars in sleeves of $1\%$, $0.3\%$ and $2.5\%$ vol is a
portfolio that is $85\%$ the loud one.
\item \textbf{Trusting the optimizer's Sharpe.} A near-zero-vol
sleeve reports an absurd Sharpe and gets 99.8\% weight ---
$1/\sigma$ is a proposal that still needs capacity and leverage
limits.
\item \textbf{Vol targeting as alpha.} Targeting controls risk,
not return --- on this tape it mechanically held vol at $44\%$
and still lost the Sharpe to buy-and-hold.
\item \textbf{Budgeting on vibes.} A drawdown rule invented
mid-cascade is a panic sale with extra steps. Commit the
thresholds before the drawdown, or do not pretend they exist.
\end{itemize}

\section*{Summary}
\begin{center}
\begin{tabular}{@{}lll@{}}
\toprule
Tool & Question answered & Failure mode \\
\midrule
Correlation sizing & is this one bet or two? & corr. estimated in calm regime \\
Vol targeting & what is today's risk? & controls vol, not return \\
Risk parity & who owns the risk? & uncapped $1/\sigma$ fantasy \\
Drawdown budget & when do we stop? & threshold invented mid-loss \\
Capacity limits & can the book hold it? & Sharpe of a drip sleeve \\
\bottomrule
\end{tabular}
\end{center}

\section*{Additional exercises}
\addcontentsline{toc}{section}{Additional exercises}

\begin{exercise}[easy]
A sleeve has 60\% annual vol and your budget is 30\%. What
vol-target weight applies, and what does the $2\times$ cap change
when realized vol falls to 15\%?
\end{exercise}
\begin{solution}
$w=0.5$ at 60\%; at 15\% the uncapped weight is $2.0$, so the
cap binds --- realized vol then exceeds target.
\end{solution}

\begin{exercise}[easy]
Two sleeves have vols 20\% and 5\%. Write the $1/\sigma$ weights
and the risk shares.
\end{exercise}
\begin{solution}
$w=(0.2,0.8)$; each sleeve contributes exactly $0.5$ of risk.
\end{solution}

\begin{exercise}[easy]
A portfolio spends 452 of 720 days under $-10\%$. What fraction
of the year is that, and why does a drawdown budget not shorten it?
\end{exercise}
\begin{solution}
$63\%$. The budget caps depth, not duration --- re-entry rules
lag the recovery, so time-under-water can lengthen.
\end{solution}

\begin{exercise}[easy]
On the pinned tape, corr(spot, funding carry) $=-0.280$. Why is
a negative correlation between a beta sleeve and a carry sleeve
plausible rather than a data bug?
\end{exercise}
\begin{solution}
Funding pays more when longs crowd (rallies), and the carry
sleeve's return is the funding stream minus small basis noise ---
its sign tracks positioning stress, not spot direction.
\end{solution}

\begin{exercise}[easy]
Your risk-parity output hands 99.8\% to a 0.5bp-vol sleeve.
Name the two constraints the formula ignored.
\end{exercise}
\begin{solution}
Leverage limits (financing the outsized book) and capacity limits
(the sleeve's liquidity cannot absorb the implied size).
\end{solution}

\begin{exercise}[medium]
Show that for two uncorrelated sleeves, $w_i\propto1/\sigma_i$
minimizes portfolio vol subject to weights summing to 1.
\end{exercise}
\begin{solution}
Minimize $w_1^2\sigma_1^2+w_2^2\sigma_2^2$ with $w_1+w_2=1$:
FOC gives $w_1\sigma_1^2=w_2\sigma_2^2$; wait --- that gives
$1/\sigma^2$. Correcting honestly: equal \emph{risk contribution}
$w_i\sigma_i$ equal gives $w\propto1/\sigma$; minimum-variance
gives $1/\sigma^2$. State which you meant --- they differ.
\end{solution}

\begin{exercise}[medium]
The vol-target rule uses trailing 20-day $\hat\sigma$. Give two
failure windows where the realized vol will overshoot the target.
\end{exercise}
\begin{solution}
Gap openings (vol arrives inside the bar, before the estimate
updates) and vol spikes faster than the 20-day window decays
(leverage stays high through the first days of a regime break).
\end{solution}

\begin{exercise}[medium]
A drawdown budget halves at $-10\%$ and exits at $-20\%$. Sketch
the equity-path where this \emph{loses} to buy-and-hold.
\end{exercise}
\begin{solution}
A $-19\%$ dip that V-shapes to new highs: the budget halves into
the bottom and re-enters late --- whipsaw cost. Budgets lose in
sharp recoveries, win in cascades.
\end{solution}

\begin{exercise}[medium]
The capstone reports carry Sharpe $+107$. Why is that number
meaningless even though it is correctly computed?
\end{exercise}
\begin{solution}
Sharpe divides by realized vol; a near-deterministic funding drip
has $\sigma\approx0.5$bp/day, so the ratio explodes. It measures
the sleeve's smoothness, not its capacity to make money at size.
\end{solution}

\begin{exercise}[medium]
Weekend vol is 25\% vs weekday 39\%. You trade both regimes at
one vol target. What systematic mistake does that bake in?
\end{exercise}
\begin{solution}
Weekends get over-levered relative to their true risk budget ---
the trailing $\hat\sigma$ lags the calendar; a regime-conditioned
target fixes it.
\end{solution}

\begin{exercise}[hard]
Derive the risk-parity weight update when sleeves are correlated:
why does $1/\sigma$ stop being the answer?
\end{exercise}
\begin{solution}
With $\Sigma\ne\mathrm{diag}$, risk contribution is
\[ \mathrm{RC}_i = w_i(\Sigma w)_i/\sigma_p ; \]
equalizing needs iterative solve (Spinu/Newton), not a closed form ---
$1/\sigma_i$ is only exact when $\Sigma$ is diagonal.
\end{solution}

\begin{exercise}[hard]
Design a funding-carry sleeve's risk budget: given the pinned
tape's 2.69bp/day mean and 0.48bp vol, what breaks the strategy
at $100\times$ scale?
\end{exercise}
\begin{solution}
Funding flips sign in stressed regimes (the sleeve becomes the
payer); perp borrow/locates and venue capacity cap the size;
cross-margin liquidation paths on the hedge leg.
\end{solution}

\begin{exercise}[hard]
Formulate the full constrained problem: maximize expected return
s.t.\ risk contributions equal, leverage $\le L$, sleeve $i$
size $\le C_i$. Which constraint binds first on our tape?
\end{exercise}
\begin{solution}
The capacity constraint on the carry sleeve --- its tiny vol
wants huge weight but real funding depth is finite; then leverage.
\end{solution}

\begin{exercise}[hard]
Your $1/\sigma$ weights were estimated on 20 days of carry data.
Estimate the standard error of $\hat\sigma$ at $n=20$ and say what
it does to the weights.
\end{exercise}
\begin{solution}
$\mathrm{se}(\hat\sigma)\approx\sigma/\sqrt{2n}\approx11\%$ ---
the dominant sleeve's weight carries $\sim$$15\%$ estimation
error; small samples invert to large weight errors.
\end{solution}

\begin{exercise}[hard]
Argue for or against: ``drawdown budgets are stop-losses for
cowards.'' Use the tape's numbers.
\end{exercise}
\begin{solution}
Against as stated: the budgeted book beat B\&H ($1.33$ vs $1.24$)
with a third the drawdown --- but the honest version is
regime-conditional: in V-shaped markets the same rule whipsaws.
The budget is insurance, and insurance has a premium.
\end{solution}

\section*{Notes}
This chapter deliberately ends on a constraint, not a formula.
Every earlier chapter asked ``is this real?''; this one asks
``how much, and who pays when it is wrong?'' --- the allocator's
question, and the one a retail book most often skips.

\section*{Further reading}
\begin{itemize}
\item T.~Roncalli, \emph{Introduction to Risk Parity and
Budgeting} --- the reference text for the $1/\sigma$ family
(borrow or library).
\item R.~Grinold \& R.~Kahn, \emph{Active Portfolio Management},
chs.~2--4 --- the institutional formulation of sizing (borrow or
library).
\item A.~Ilmanen, \emph{Expected Returns} --- the allocator's
survey of what drives sleeve returns (borrow or library).
\item Risk parity, overview with further links:

\url{https://en.wikipedia.org/wiki/Risk_parity}
\item Portfolio construction theory:

\url{https://en.wikipedia.org/wiki/Modern_portfolio_theory}
\item Position sizing via the growth-optimal rule:

\url{https://en.wikipedia.org/wiki/Kelly_criterion}
\item Portfolio Visualizer --- free backtester for sanity-checking
allocations:

\url{https://www.portfoliovisualizer.com}
\end{itemize}

\begin{rulebox}[Watch instead]
\begin{itemize}
\item ``How to Build a Recession-Proof Investment Portfolio'' ---
Real Vision, Chris Cole interviewed by D.~DiMartino Booth
(52:39). Vol targeting, risk parity, and crisis-era allocation
by a practitioner; opened and verified 2026-10-11.
\url{https://www.youtube.com/watch?v=SkfgEZtJ9LA}
\item ``How The Economic Machine Works'' --- Ray Dalio,
Principles (31:00). The debt-cycle macro backdrop allocators
size against; opened and verified 2026-10-11.
\url{https://www.youtube.com/watch?v=PHe0bXAIuk0}
\end{itemize}
\end{rulebox}
